% =========================================================
% Bounds, Suprema, Infima + Table + Figure
% =========================================================

\section{Bounds and Extremal Values}

\begin{remark*}[Predicate vs.\ Existence: Structural Template]
Each bound-type concept (maximum, minimum, upper bound, supremum, etc.) is
defined by a predicate relative to an ambient ordered set $S$. Given
$A \subseteq S$, let $P_A(x)$ denote such a predicate.

\medskip
\noindent\textbf{Predicate form.}
The statement ``$x$ is a $P$-object of $A$'' means that $P_A(x)$ holds for a
specific identified element $x \in S$. To prove this, one exhibits the element
and verifies each condition in $P_A(x)$.

\medskip
\noindent\textbf{Existence form.}
The statement ``$A$ has a $P$-object'' means that
\[
\exists x \in S \; P_A(x).
\]
Once established, one may fix such an $x$ and reason with the fact that
$P_A(x)$ holds.

\medskip
\noindent\textbf{Notational glosses.}
\begin{itemize}
  \item ``$x$ is a $P$-object of $A$'' abbreviates $P_A(x)$.
  \item ``$A$ has a $P$-object'' abbreviates $\exists x \in S \; P_A(x)$.
\end{itemize}

\medskip
\noindent\textbf{Proof strategy summary.}

\begin{center}
\renewcommand{\arraystretch}{1.4}
\begin{tabular}{@{} p{4cm} p{3cm} p{6cm} @{}}
\hline
\textbf{Goal} & \textbf{Form} & \textbf{Strategy} \\
\hline
$P_A(x)$ & Predicate & Exhibit $x$ and verify $P_A(x)$ directly. \\
\noalign{\vskip 0.3em}
$\exists x \in S \; P_A(x)$ & Existence & Construct or identify a candidate and verify $P_A(x)$. \\
\noalign{\vskip 0.3em}
\textit{Hypothesis:} $\exists x \in S \; P_A(x)$ & --- & Fix such an $x$ and use the fact that $P_A(x)$ holds. \\
\hline
\end{tabular}
\end{center}

\noindent\emph{Note.}
The hypothesis row is not a proof goal. It records the standard move of fixing
an existential witness once existence has been established.
\end{remark*}

\begin{center}
\input{volume-iii/analysis/bounding/notes/bounds-extremals/figure-bounds}
\end{center}

\begin{center}
\small
\textbf{Bounds, Extrema, Infimum, and Supremum for Subsets of $S$}
\end{center}

\medskip

\begin{center}
\renewcommand{\arraystretch}{1.6}
\begin{tabular}{@{} p{2.4cm} p{6.2cm} p{3.0cm} p{3.2cm} @{}}
\hline
\textbf{Concept} & \textbf{Predicate $P_A(x)$} & \textbf{Predicate Form} & \textbf{Existence Form} \\
\hline
Upper bound
&
$\forall a \in A \; (a \le x)$
&
$x$ is an upper bound of $A$
&
$A$ is bounded above$
$
\\[0.2em]

Lower bound
&
$\forall a \in A \; (x \le a)$
&
$x$ is a lower bound of $A$
&
$A$ is bounded below
$
$
\\
\hline

Maximum
&
$(x \in A) \land \forall a \in A \; (a \le x)$
&
$x$ is a maximum of $A$
&
$A$ has a maximum
$
$
\\[0.2em]

Minimum
&
$(x \in A) \land \forall a \in A \; (x \le a)$
&
$x$ is a minimum of $A$
&
$A$ has a minimum
$
$
\\
\hline

Supremum
&
$UpperBound(x,A)
\land
\forall u \in S \; (UpperBound(u,A) \rightarrow x \le u)$
&
$x$ is the supremum of $A$
&
$A$ has a supremum
$
$
\\[0.2em]

Infimum
&
$LowerBound(x,A)
\land
\forall \ell \in S \; (LowerBound(\ell,A) \rightarrow \ell \le x)$
&
$x$ is the infimum of $A$
&
$A$ has an infimum
$
$
\\
\hline
\end{tabular}
\end{center}

\medskip

\begin{remark*}[Reading the summary table]
The predicate column gives the condition that must hold for a specific
identified element $x \in S$. The existence column records the corresponding
existential statement. The rows are grouped by logical complexity: bounds use a
single universal condition, extrema add membership, and suprema and infima add
the least or greatest comparison against all competing bounds.
\end{remark*}












% =========================================================
% Bounding Definitions
% House format
% =========================================================

\subsection*{Bounds and Extremal Elements}

\begin{tcolorbox}[colback=propbox, colframe=propborder, arc=2pt,
  left=6pt, right=6pt, top=4pt, bottom=4pt,
  title={\small\textbf{Definition (Lower Bound)}},
  fonttitle=\small\bfseries]
\begin{definition}
\label{def:lower-bound}
Let $A \subseteq S$ and let $l \in S$. The element $l$ is called a
\emph{lower bound} of $A$ if
\[
l \le x
\qquad \text{for every } x \in A.
\]
\end{definition}
\end{tcolorbox}

\begin{remark*}[Standard quantified statement]
\[
l \text{ is a lower bound of } A
\quad\Longleftrightarrow\quad
\forall x \in A \; (l \le x).
\]
\end{remark*}

\begin{remark*}[Predicate statement]
\[
LowerBound(l,A).
\]
\end{remark*}

\begin{remark*}[Negated quantified statement]
\[
l \text{ is not a lower bound of } A
\quad\Longleftrightarrow\quad
\exists x \in A \; (x < l).
\]
This means there exists an element of $A$ lying strictly below $l$, so $l$ fails to bound $A$
from below.
\end{remark*}

\begin{tcolorbox}[colback=propbox, colframe=propborder, arc=2pt,
  left=6pt, right=6pt, top=4pt, bottom=4pt,
  title={\small\textbf{Definition (Upper Bound)}},
  fonttitle=\small\bfseries]
\begin{definition}
\label{def:upper-bound}
Let $A \subseteq S$ and let $u \in S$. The element $u$ is called an
\emph{upper bound} of $A$ if
\[
x \le u
\qquad \text{for every } x \in A.
\]
\end{definition}
\end{tcolorbox}

\begin{remark*}[Standard quantified statement]
\[
u \text{ is an upper bound of } A
\quad\Longleftrightarrow\quad
\forall x \in A \; (x \le u).
\]
\end{remark*}

\begin{remark*}[Predicate statement]
\[
UpperBound(u,A).
\]
\end{remark*}

\begin{remark*}[Negated quantified statement]
\[
u \text{ is not an upper bound of } A
\quad\Longleftrightarrow\quad
\exists x \in A \; (u < x).
\]
This means there exists an element of $A$ lying strictly above $u$, so $u$ fails to bound $A$
from above.
\end{remark*}

\begin{tcolorbox}[colback=propbox, colframe=propborder, arc=2pt,
  left=6pt, right=6pt, top=4pt, bottom=4pt,
  title={\small\textbf{Definition (Minimum)}},
  fonttitle=\small\bfseries]
\begin{definition}
\label{def:minimum}
Let $A \subseteq S$ and let $m \in S$. The element $m$ is called a
\emph{minimum} of $A$, denoted $m = \min A$, if $m \in A$ and
\[
m \le x
\qquad \text{for every } x \in A.
\]
\end{definition}
\end{tcolorbox}

\begin{remark*}[Standard quantified statement]
\[
m = \min A
\quad\Longleftrightarrow\quad
(m \in A) \land \forall x \in A \; (m \le x).
\]
\end{remark*}

\begin{remark*}[Predicate statement]
\[
Minimum(m,A).
\]
\end{remark*}

\begin{remark*}[Negated quantified statement]
\[
m \neq \min A
\quad\Longleftrightarrow\quad
(m \notin A) \lor \exists x \in A \; (x < m).
\]
This means either $m$ is not an element of $A$, or else some element of $A$ lies strictly
below $m$, so $m$ cannot be the least element of $A$.
\end{remark*}

\begin{tcolorbox}[colback=propbox, colframe=propborder, arc=2pt,
  left=6pt, right=6pt, top=4pt, bottom=4pt,
  title={\small\textbf{Definition (Maximum)}},
  fonttitle=\small\bfseries]
\begin{definition}
\label{def:maximum}
Let $A \subseteq S$ and let $m \in S$. The element $m$ is called a
\emph{maximum} of $A$, denoted $m = \max A$, if $m \in A$ and
\[
x \le m
\qquad \text{for every } x \in A.
\]
\end{definition}
\end{tcolorbox}

\begin{remark*}[Standard quantified statement]
\[
m = \max A
\quad\Longleftrightarrow\quad
(m \in A) \land \forall x \in A \; (x \le m).
\]
\end{remark*}

\begin{remark*}[Predicate statement]
\[
Maximum(m,A).
\]
\end{remark*}

\begin{remark*}[Negated quantified statement]
\[
m \neq \max A
\quad\Longleftrightarrow\quad
(m \notin A) \lor \exists x \in A \; (m < x).
\]
This means either $m$ is not an element of $A$, or else some element of $A$ lies strictly
above $m$, so $m$ cannot be the greatest element of $A$.
\end{remark*}









\begin{tcolorbox}[colback=propbox, colframe=propborder, arc=2pt,
  left=6pt, right=6pt, top=4pt, bottom=4pt,
  title={\small\textbf{Definition (Supremum)}},
  fonttitle=\small\bfseries]
\begin{definition}
\label{def:supremum}
Let $A \subseteq S$ and let $s \in S$. The element $s$ is called the
\emph{supremum} of $A$, denoted $s = \sup A$, if $s$ is an upper bound of $A$
and $s$ is less than or equal to every upper bound of $A$.
\end{definition}
\end{tcolorbox}

\begin{remark*}[Standard quantified statement]
\[
s = \sup A
\quad\Longleftrightarrow\quad
\bigl(\forall x \in A \; (x \le s)\bigr)
\land
\bigl(\forall u \in S \; ((\forall x \in A \; (x \le u)) \rightarrow s \le u)\bigr).
\]
\end{remark*}

\begin{remark*}[Definition predicate reading]
\[
Supremum(s,A)
\quad\Longleftrightarrow\quad
\underbrace{UpperBound(s,A)}_{\text{$s$ is an upper bound of $A$}}
\land
\underbrace{\forall u \in S \; (UpperBound(u,A) \rightarrow s \le u)}_{\text{$s$ lies below every upper bound of $A$}}.
\]
\end{remark*}

\begin{remark*}[Negated quantified statement]
\[
s \neq \sup A
\quad\Longleftrightarrow\quad
\bigl(\exists x \in A \; (s < x)\bigr)
\lor
\bigl(\exists u \in S \; ((\forall x \in A \; (x \le u)) \land u < s)\bigr).
\]
\end{remark*}

\begin{remark*}[Failure mode decomposition]
\[
s \ne \sup A
\quad\Longleftrightarrow\quad
\underbrace{\exists x \in A \; (s < x)}_{\text{$s$ fails to be an upper bound}}
\;\lor\;
\underbrace{\exists u \in S \; \bigl((\forall x \in A \; (x \le u)) \land u < s\bigr)}_{\text{$s$ fails to be the least upper bound}}.
\]
\end{remark*}

\begin{remark*}[Negation predicate reading]
\[
\neg Supremum(s,A)
\quad\Longleftrightarrow\quad
\underbrace{\neg UpperBound(s,A)}_{\text{$s$ is not an upper bound}}
\lor
\underbrace{\exists u \in S \; (UpperBound(u,A) \land u < s)}_{\text{there is a smaller upper bound}}.
\]
\end{remark*}

\begin{remark*}[Failure modes]\hfill
\begin{itemize}
  \item \textbf{Failure to be an upper bound.} There exists an element of $A$ strictly larger than $s$, so $s$ does not bound $A$ from above at all.
  \item \textbf{Failure to be the least upper bound.} The element $s$ may still be an upper bound, but there exists a smaller upper bound $u < s$, so $s$ is not the least one.
\end{itemize}
\end{remark*}














\begin{tcolorbox}[colback=propbox, colframe=propborder, arc=2pt,
  left=6pt, right=6pt, top=4pt, bottom=4pt,
  title={\small\textbf{Definition (Infimum)}},
  fonttitle=\small\bfseries]
\begin{definition}
\label{def:infimum}
Let $A \subseteq S$ and let $t \in S$. The element $t$ is called the
\emph{infimum} of $A$, denoted $t = \inf A$, if $t$ is a lower bound of $A$
and every lower bound of $A$ is less than or equal to $t$.
\end{definition}
\end{tcolorbox}

\begin{remark*}[Standard quantified statement]
\[
t = \inf A
\quad\Longleftrightarrow\quad
\bigl(\forall x \in A \; (t \le x)\bigr)
\land
\bigl(\forall l \in S \; ((\forall x \in A \; (l \le x)) \rightarrow l \le t)\bigr).
\]
\end{remark*}

\begin{remark*}[Predicate statement]
\[
Infimum(t,A).
\]
Equivalently,
\[
LowerBound(t,A)
\land
\forall l \in S \; (LowerBound(l,A) \rightarrow l \le t).
\]
\end{remark*}

\begin{remark*}[Negated quantified statement]
\[
t \neq \inf A
\quad\Longleftrightarrow\quad
\bigl(\exists x \in A \; (x < t)\bigr)
\lor
\bigl(\exists l \in S \; ((\forall x \in A \; (l \le x)) \land t < l)\bigr).
\]
This means either $t$ fails to be a lower bound of $A$, or $t$ is a lower bound but not the
greatest one.
\end{remark*}










% ---------------------------------------------------------
\subsection{Consequences}

\begin{remark*}[Logical structure]
Bounds, extrema, and least/greatest bounds form a hierarchy of increasing
logical complexity.
Upper and lower bounds are defined by a universal quantifier alone.
Maximum and minimum add a membership condition: the extremum must belong
to the set.
Supremum and infimum relax membership but add a minimality (or maximality)
condition: the bound must be the \emph{least} (or \emph{greatest}) among
all bounds.
Maximum implies supremum --- when $\max A$ exists, it equals $\sup A$
(Proposition~\ref{prop:sup-unique} and the proof at
\hyperref[prf:max-is-sup]{max-is-sup}) --- but the converse fails:
$(0,1)$ has $\sup = 1$ but no maximum.
\end{remark*}

\begin{remark*}[Connections forward]
The $\varepsilon$-characterization (Proposition~\ref{prop:eps-char}) is the
workhorse tool of the bounding chapter.
Every proof in the Bound Arithmetic section
(\S\ref{sec:bound-arithmetic}) invokes it to verify that a candidate value
is the least upper bound, not merely an upper bound.
The Monotone Convergence Theorem (convergence chapter) uses the same
mechanism: the limit of a bounded monotone sequence is identified as
$\sup\{a_n\}$, and the $\varepsilon$-characterization produces the required
$N$.
\end{remark*}



